%%%%%
%
\section{FSSP: Correction des co\"{\i}ncidences avec le
taux de comptage}

\itemstyle{Op\'erateur:} {\sf fssp\_ncor\_specrs}
\index{fssp\_ncor\_specrs}

%%%%%OBJET%%%%
\itemstyle{Objet:}
Calcul du nombre de goutelettes total en corrig\'e des co\"{\i}ncidences; m\'ethode
utilisant uniquement le taux de comptage.

%%%%%ARGUMENTS%%%%
\bigskip
%%%%%ENTREES%%%%
\begin{tabarg}
\multicolumn{3}{l}{\titretab {Entr\'ees}}\\
\hline
$Strobe$ & Vecteur, {\tt float} & Nombre de particules dans la D.O.F.
   (comptes par p\'eriode d'\'echantillonnage)\\
$Reset$& Vecteur, {\tt float} & Nombre total de particules compt\'ees 
    (comptes par p\'eriode d'\'echantillonnage)\\
$\tau_{spec}$ & Vecteur, {\tt float} & temps de transit estim\'e \`a
          partir du spectre de gouttes ($s$)\\
$f_e$ & Coeff., {\tt float} & Fr\'equence d'\'echantillonnage ($Hz$)\\
$\tau_1$ & Coeff., {\tt float} & D\'elai \'electronique court ($s$)\\
$\tau_2$ & Coeff., {\tt float} & D\'elai \'electronique long ($s$)\\
$\tau_r$ & Coeff., {\tt float} & D\'elai \'electronique pour le reset ($s$)\\
$Crt_0$ & Coeff., {\tt float} & Taux de comptage minimal au dessus
duquel on peut effectuer la correction ($count.second^{-1}$)\\
\end{tabarg}

\medskip
%%%%%SORTIES%%%%
\begin{tabarg}
\multicolumn{3}{l}{\titretab {Sorties}}\\
\hline
$N_{rs}$ & Vecteur, {\tt float} & Nombre total de gouttelettes
corrig\'e des co\"{\i}ncidences par le taux de comptage
(\textit{comptes par p\'eriode d'\'echantillonnage}) \\
\end{tabarg}


%%%%%FORMULE%%%%
\itemstyle{Formule}

%%fssp_ncor_specrs

L'algorithme et les formules \'etant complexes
seul le principe de calcul est pr\'esent\'e ici.\\
Selon Brenguier-89 \cite{Brenguier-89}, 
le nombre de gouttes mesur\'ees par un FSSP
s'exprime comme une fonction~: \\ 
\begin{eqnarray}
N & = & f_n ( N_c, \tau, \tau_1, \tau_2, \tau_r ) \label{eqn:crt2}
\end{eqnarray}
o\`u~:\\
$N$ : nombre total de gouttes compt\'ees \\
$N_c$ : nombre total de gouttes corrig\'e des
co\"{\i}ncidences \\
$\tau $ : temps de transit d'une particule \\
\\
On pose $\tau = \tau_{spec}$, et $N=Reset$.
Ainsi, cette \'equation ne
poss\`ede plus qu'une inconnue $N_c$. 
Elle est r\'esolue par it\'erations successives.   
(\ref{eqn:crt2}) donne $N_{rs}$.\\
Si $Reset\times f_e \le Crt_0$ alors aucune correction
est calcul\'ee et $N_{ac}=0$ et $N_{rs}=0$.


%%%%%AUTEUR%%%%
\itemstyle{Commentaire:}
Le calcul est identique \`a celui fait dans \textsf{fssp\_ncor\_specacr},
sans prendre en compte l'activit\'e. Le r\'esultat est de moindre
qualit\'e, on n'utilisera donc cette formule que dans les cas o\`u
l'activit\'e n'a pas \'et\'e mesur\'ee.

%%%%%AUTEUR%%%%
\itemstyle{Auteur:}
CNRM/GMEI/TRAMM

%%%%%REFERENCE%%%%
\itemstyle{R\'ef\'erences:}
Brenguier-89\cite{Brenguier-89}
